% Source: https://artofproblemsolving.com/community/c3h2417380_polynomials
% Stable identity: aops-topic-2417380-q01
The right side is a perfect square:
\[\left(x+\dfrac{1}{2x}\right)^2=x^2+2\cdot x\cdot\dfrac{1}{2x}+\dfrac{1}{4x^2}=x^2+1+\dfrac{1}{4x^2},\qquad\text{and}\qquad x+\dfrac{1}{2x}=\dfrac12\left(2x+\dfrac1x\right).\]
So, writing $t=2x+\dfrac1x$, the condition says $f(t)=\dfrac{t^2}{4}$ for every number $t$ of this form.

Every $t\ge2\sqrt2$ is of this form: $2x+\dfrac1x=t$ means $2x^2-tx+1=0$, whose discriminant $t^2-8$ is nonnegative, and whose roots have positive product $\dfrac12$ and positive sum $\dfrac t2$, so they are positive real numbers. (By AM-GM, $2x+\dfrac1x\ge2\sqrt2$ for $x>0$, so these are exactly the values for $x>0$.)

Hence the polynomial $f(t)-\dfrac{t^2}{4}$ has infinitely many roots, so it is the zero polynomial, and
\[\boxed{f(x)=\dfrac{x^2}{4}}.\]
(Without the assumption that $f$ is a polynomial, the condition would fix $f(t)$ only for $|t|\ge2\sqrt2$, the values taken by $2x+\dfrac1x$.)
